\subsection{First-Order Theories with Equality}

In this chapter we introduce the standard interpretation, i.e. a first-order theory with equality. Let $A_1^2$ be a predicate letter for which $A_1^2(t,s)$ indicates that $t=s$, and vice versa $\lnot A_1^2(t, s)$ indicates that $t \not = s$. More specifically, K is a first-order theory of equality if the followin are theorems of K:

\begin{enumerate} [label = (A\arabic*)]
    \setcounter{enumi}{5}
    \item
        \begin{tabular}{p{8cm} r}
            $(\forall x_1)x_1=x_1$ & (reflexivity principle)
        \end{tabular}
    \item
        \begin{tabular}{p{8cm} r}
            $x=y \imp \big(\B(x,x) \imp \B(x,y) \big)$ & (substitutivity of equality)
        \end{tabular}
\end{enumerate}
Where $x$ and $y$ are any variables, $\B(x,y)$ is any wf, and $\B(x,y)$ arises from $\B(x,x)$ by replacing some (but not necessarily all) free occurrences of $x$ by $y$ provided that $y$ is free for $x$ in $\B(x,x)$. So, $\B(x,y)$ may or may not contain free occurrences of $x$. Note that the numbering of the above axioms is a continuation of the logical axioms (A1)-(A5).

\textbf{Proposition 2.23}

In any theory of equality,
\begin{enumerate}
    \item $\vdash t=t$ for any term $t$.
    \item $\vdash t=s \imp s=t$ for any terms $t$ and $s$.
    \item $\vdash t=s \imp (s=r \imp t=r)$ for any terms $t$, $s$, and $r$.
\end{enumerate}

\textbf{Proof}
\begin{enumerate} [label = \alph*.]
    \item By (A6), $\vdash (\forall x_1) x_1=x_1$. Hence, by rule A4, $\vdash t=t$.
    \item \begin{enumerate} [label = \roman*.]
        \item Let $x$ and $y$ be variables not occurring in $t$ or $s$.
        \item Letting $\B(x,x)$ be $x=x$ and $\B(x,y)$ be $x=y$ in schema (A7): \[\vdash x=y \imp (x=x \imp x=y)\]
        \item But, by (a), $\vdash x=x$.
        \item So, by instance of the tautology
        \[(A\imp (B \imp C)) \imp (B \imp (A \imp C))\]
        and two applications of MP, we get $\vdash x=y \imp y=x$.
        \item Two applications of Gen yields $\vdash(\forall x)(\forall y)(x=y \imp y=x)$.
        \item Then two applications of rule A4 yields $\vdash t=s \imp s=t$.
    \end{enumerate}
    \item \begin{enumerate} [label = \roman*.]
        \item Let $x$, $y$, and $z$ be three variables not occurring in $t$, $s$, or $r$.
        \item Then let $\B(y,y)$ be $y=z$ and $\B(x,x)$ be $x=z$ in (A7).
        \item Now, with $x$ and $y$ interchanged, we obtain $\vdash y=x \imp(y=z \imp x=z)$.
        \item By, by (b), $\vdash x=y \imp y=x$.
        \item Now, using an instance of the tautology
        \[(A\imp B) \imp ((B \imp C) \imp (A \imp C))\]
        and two applications of MP, we obtain $\vdash x=y\imp(y=z\imp x=z)$.
        \item Similarly to our proof of (b), by three applications of gen we obtain 
        \[(\forall x)(\forall y)(\forall z)(x=y \imp (y=z \imp x=z))\]
        \item Then three applications of rule A4 obtains $\vdash t=s \imp (s=r \imp t=r)$.
    \end{enumerate}
\end{enumerate}

\textbf{Proposition 2.24}
Let K be a theory for which (A6) holds and (A7) holds for all atomic wfs $\B(x,x)$ in which there are no individual constants. Then K is a theory with equality, that is, (A7) holds for all wfs $\B(x,x)$.


\textbf{Proof}
This proposition serves the purpose of reducing the schema (A7) to simpler cases. I don't feel like writing this proof up at the moment and it doesn't seem incredibly significant so I am skipping it for now.

\textbf{Proposition 2.25}
Let K be any theory in which (A6) holds and the following are true:
\begin{enumerate} [label = \alph*.]
    \item Schema (A7) holds for all atomic wfs $\B(x,x)$ such that no function letters or individual constants occur in $\B(x,x)$ and $\B(x,y)$ comes from $\B(x,x)$ by replacing exactly one occurrence of $x$ by $y$.
    \item $\vdash x=y \imp f_j^n(z_1,\dots,z_n)=f_j^n(w_1,\dots,w_n)$, where $f_j^n$ is any function letter of K, $z_1,\dots, z_n$ are variables, and $f_j^n(w_1,\dots,w_n)$ arises from $f_j^n(z_1,\dots,z_n)$ by replacing exactly one occurrence of $x$ by $y$.
\end{enumerate}
Then K is a theory of equality. 

\textbf{Proof}
\begin{enumerate}
    \item By repeated application, our assumptions can be extended to replacements of ore than one occurrence of $x$ by $y$.
    \item Also, Poposition 2.23 is still derivable.
    \item By Proposition 2.24, it suffices to prove (A7) for only atomic wfs without individual constants.
    \item But, hypothesis (a) allows us to prove 
    \[
    \vdash (y_1=z_1 \land \dots \land y_n=z_n) \imp (\B(y_1, \dots, y_n) \imp \B(z_1, \dots, z_n))
    \]
    for all variables $y_1, \dots, y_n, z_1, \dots z_n$ and any atomic wf $\B(y_1, \dots, y_n)$ without function letters or individual constants.
    \item So, if suffices to show the following:
    \begin{center}
        (*) If $t(x,x)$ is a term without individual constants and $t(x,y)$ comes from $t(x,x)$ by replacing some occurrences of $x$ by $y$, then $\vdash x=y \imp t(x,x)=t(x,y)$. \footnote{$\ $You can clarify how (*) is applied by using it to prove the following instance of (A7): $\vdash x=y \imp (A_1^1(f_1^1(x))\imp A_1^1(f_1^1(y)))$. Let $t(x,x)$ be $f_1^1(x)$ and let $t(x,y)$ by $f_1^1(y)$.}
    \end{center}
    \item But (*) can be proved using hypothesis (b) by induction on the number of function letters in $t(x,x)$. This is left as an exercise in the book, maybe I will add this later (probably not, sorry).
    \item It can be seen from Proposition 2.25 that it is only necessary to verify (A7) for a finite list of special cases when the language K has only finitely many predicate and function letters. \footnote{$\ $In fact, $n$ wfs for each $A_j^n$ and $n$ wfs for each $f_j^n$.}
\end{enumerate}

Some Additional notes. Not sure why this isn't a proposition, considering that this part of the textbook seems very much like a proof.

\begin{enumerate}
    \item In any model for a theory K with equality, the relation E in the model corresponding to the predicate letter = is an equivalence relation (by Proposition 2.23).
    \item If this relation E is the identity relation in the domain of the model, then the model is said to be \textit{normal}.
    \item Any model M for K can be \textit{contracted} to a normal model M$^*$ for K by taking the domain $D^*$ of M$^*$ to be the set of equivalence classes determined by the relation E in the domain $D$ of M.
    \item For a predicate letter $A_j^n$ and for any equivalence classes $[b_1], \dots, [b_n]$ in $D^*$ determined by elements $b_1, \dots, b_n$ in $D$, we let $(A_j^n)^{\text{M}^*}$ hold for $([b_1], \dots, [b_n])$ iff $\model{A_j^n}$ holds for $(b_1, \dots, b_n)$.
    \item Notice that it makes no difference which representatives $b_1, \dots, b_n$ we select in the given equivalence classes because, from (A7):
    \[
    \vdash x_1=y_1 \land \dots \land x_n = y_n \imp (A_j^n(x_1, \dots, x_n) \bicon A_j^n(y_1, \dots, y_n))
    \]
    \item Likewise, for any function letter $f_j^n$ and any equivalence classes $[b_1], \dots, [b_n]$ in $D^*$, let $(f_j^n)^{\text{M}^*}([b_1], \dots, [b_n])=[\model{f_j^n}(b_1, \dots, b_n)]$.
    \item Again note that this is independent of the choise of the representatives $b_1, \dots, b_n$, since, from (A7), we can prove:
    \[
    \vdash x_1=y_1 \land \dots \land x_n = y_n \imp (f_j^n(x_1, \dots, x_n) \bicon f_j^n(y_1, \dots, y_n))
    \]
    \item For any individual constant $a_i$ let $(a_i)^{\text{M}^*}=[\model{a_i}]$.
    \item The relation E$^*$ corresponding to = in the model M$^*$ is the identity relation in $D^*:E^*([b_1],[b_2])$ iff $E(b_1,b_2)$, that is, iff $[b_1]=[b_2]$.
    \item Now one can prove by induction the following lemma:
    \item If $s=(b_1, b_2, \dots)$ is a denumerable sequence of equivalence classes, then a wf $\B$ is satisfied by $s$ in M iff $\B$ is satisfied by $s'$ in M$^*$.
    \item It follows that, for any wf $\B$, $\B$ is true for M iff $\B$ is true for M$^*$.
    \item Hence, because M is a model of K, M$^*$ is a normal model of K.
\end{enumerate} \\


\textbf{Proposition 2.26 (Extension of Proposition 2.17)}
(\godel, 1930) Any consistent theory with equality K has a finite or denumerable normal model.

\textbf{Proof}

By Proposition 2.17, K has a denumerable model M. Hence, the contraction of M to a normal model yields a finite or denumerable normal model M$^*$ because the set of equivalence classes in a denumerable set $D$ is either finite or denumerable.

\textbf{Corollary 2.27 (Extension of Skolem-L\"owenheim Theorem)}
Any theory with equality K that has an infinite normal model M has a denumerable normal model.

\textbf{Proof}
\begin{enumerate}
    \item Add to K the denumerably many new individual constants $b_1, b_2, \dots$ together with the axioms $b_i \not = b_j$ for $i \not = j$
    \item Then the new theory K$'$ is consistent.
    \item If K$'$ were inconsistent, there would be a proof in K$'$ of a contradiction $\C \land \lnot \C$, where we may assume that $\C$ is a wf of K. \footnote{$\ $If K$'$ is inconsistent, it can prove anything, including a wf $\C$ of K and its negation $\lnot \C$.}
    \item But this proof uses only a finite number of the new axioms: $b_{i1} \not = b_{j1}, \dots, b_{in} \not = b_{jn}$.
    \item Now, M can be extended to a model M$^\#$ of K plus the axioms $b_{i1} \not = b_{j1}, \dots, b_{in} \not = b_{jn}$.
    \item In fact, since M is an infinite normal model, we can choose interpretations of $b_{i1}, b_{j1}, \dots, b_{in}, b_{jn}$ so that the wfs $b_{i1} \not = b_{j1}, \dots, b_{in} \not = b_{jn}$ are true.
    \item But, since $\C \land \lnot \C$ is derivable from these wfs and the axioms of K, it would follow that $\C \land \lnot \C$ is true for M$^\#$, which is impossible.
    \item Hence, K$'$ must be consistent.
    \item Now, by Proposition 2.26, K$'$ has a finite or denumerable normal model N.
    \item But, since, for $i \not = j$, the wfs $b_i \not = b_j$ are axioms of K$'$, they are true for N.
    \item Thus, the elements in the domain of N that are interpretations of $b_1, b_2, \dots$ must be distinct, which implies that the domain of N is infinite, and therefore, denumerable.
\end{enumerate}
